% ============================================================
% The Three-Tier Theorem (proved April 2022)
% ============================================================

\documentclass[12pt]{amsart}

\usepackage{amsmath,amssymb,amsthm}
\usepackage{booktabs}
\usepackage{microtype}
\usepackage{url}

\newtheorem{theorem}{Theorem}
\newtheorem{lemma}[theorem]{Lemma}
\newtheorem{proposition}[theorem]{Proposition}
\newtheorem{corollary}[theorem]{Corollary}
\newtheorem{conjecture}[theorem]{Conjecture}
\theoremstyle{definition}
\newtheorem{definition}[theorem]{Definition}
\theoremstyle{remark}
\newtheorem{remark}[theorem]{Remark}

\title{The Three-Tier Theorem}

\author{Alexander S.\ Petty}

\email{alexander.petty@gmail.com}
\date{February 2022 (revised April 2026)}
\subjclass[2020]{11A63, 11B83}

\begin{document}
\begin{abstract}
This note completes the three-tier classification of
positive integers by the golden threshold $1/\varphi$.
Tiers~1 and~2 (smooth and golden-family integers) are
established in the companion alignment
papers~\cite{paper1,paper2}. Here we prove Tier~3:
every integer outside these families has
$\alpha(n) < 1/\varphi$.

The proof rests on the nines complement: non-terminating
fractions $k/n$ and $(n{-}k)/n$ have complementary digits
and cannot both match $1/n$ at any position. This bounds
the average non-terminating alignment to~$1/2$, giving
$\alpha(n) \le (t{+}1)/(2t)$ for rough part~$t$. Since
$4/7 < 1/\varphi$ and the smallest admissible rough part
$\ge 4$ in base~$10$ is $t = 7$, all non-golden integers
fall below the threshold.
\end{abstract}
\maketitle

% ============================================================
\section{Introduction}

This paper proves the three-tier classification.

The first paper~\cite{paper1} introduced repetend
alignment and identified the golden threshold
$1/\varphi$. Papers~2 and~3~\cite{paper2,paper3}
established the alignment formula for all primes, with
the golden bound $\alpha(pm) < 1/\varphi$ for $p \ge 5$.
The classification is:

\begin{enumerate}
\item[\textbf{Tier~1.}] $n$ is $10$-smooth: $\alpha(n) = 1$.
\item[\textbf{Tier~2.}] $n = 3m$, $m$ is $10$-smooth, $m \ge 4$:
      $1/\varphi \le \alpha(n) < 1$.
\item[\textbf{Tier~3.}] All other $n$: $\alpha(n) < 1/\varphi$.
\end{enumerate}

Tiers~1 and~2 are proved in~\cite{paper1}. Tier~3 is
proved here by the nines complement bound: for any
$n$ with rough part $t \ge 7$ (coprime to the base),
$\alpha(n) \le (t{+}1)/(2t) \le 4/7 < 1/\varphi$.
The cases $t = 3$ with small smooth part are verified
directly.

% ============================================================
\section{The Rough Part}

\begin{definition}
For a positive integer $n$ and a base $b$, the \emph{smooth part}
$s(n)$ is the largest divisor of $n$ whose prime factors all
divide~$b$. The \emph{rough part} is $t(n) = n / s(n)$.
\end{definition}

In base~$10$, the smooth part absorbs all factors of $2$ and $5$,
and the rough part collects all other prime factors. For example:
$s(60) = 20$, $t(60) = 3$; and $s(28) = 4$, $t(28) = 7$.

\begin{lemma}\label{lem:terminating-count}
The number of terminating fractions in $\{k/n : 1 \le k \le n{-}1\}$
is $s(n) - 1$.
\end{lemma}

\begin{proof}
The fraction $k/n$ terminates if and only if $n / \gcd(k,n)$ is
$b$-smooth, which occurs if and only if $t(n)$ divides~$k$. The
multiples of $t(n)$ in $\{1, \ldots, n{-}1\}$ number
$n/t(n) - 1 = s(n) - 1$.
\end{proof}

% ============================================================
\section{The Nines Complement Bound}

\begin{lemma}[Complement pairing]\label{lem:complement}
For $n = tm$ with rough part $t$ coprime to the even
base~$b$: if $k$ is non-terminating ($t \nmid k$), then
$k/n$ and $(n{-}k)/n$ are digit complements in the
periodic part of their expansions. At each repetend
position, their periodic digits sum to $b - 1$. In
particular:
\[
a(k/n,\, 1/n) + a((n{-}k)/n,\, 1/n) \le 1.
\]
\end{lemma}

\begin{proof}
After the transient prefix (determined by the smooth
part~$m$), the periodic digits of $k/n$ are governed by
remainders modulo~$t$. For a remainder $q$ with
$1 \le q \le t{-}1$: $\gcd(t, b) = 1$ ensures
$bq/t \notin \mathbb{Z}$, so
\[
\left\lfloor \frac{b(t{-}q)}{t} \right\rfloor
= b - \left\lceil \frac{bq}{t} \right\rceil
= b - 1 - \left\lfloor \frac{bq}{t} \right\rfloor.
\]
The periodic digits of $k/n$ and $(n{-}k)/n$ therefore
sum to $b - 1$ at each position.

For even~$b$: $(b{-}1)/2$ is not an integer, so a digit
$d$ and its complement $b{-}1{-}d$ are always distinct.
Two complementary fractions cannot both match $1/n$ at
the same periodic position.
\end{proof}

\begin{lemma}\label{lem:upper-bound}
For $n = tm$ with rough part $t$ and smooth part $m$:
\[
\alpha(n) \le \frac{m(t{+}1)/2 - 1}{tm - 1}.
\]
For $m \ge 1$ this is at most $(t{+}1)/(2t)$.
\end{lemma}

\begin{proof}
The $m - 1$ terminating fractions each contribute
alignment~$1$. The $(t{-}1)m$ non-terminating fractions
pair into complement pairs $(k, n{-}k)$ (both are
non-terminating since $t \mid k$ iff $t \mid n{-}k$).
The only possible fixed point $k = n/2$ would require
$t \mid n/2 = tm/2$; since $t$ is odd (coprime to
the even base), this means $2 \mid m$, so $n/2$ is a
multiple of~$t$ and hence terminating. No non-terminating
fraction is unpaired.
By Lemma~\ref{lem:complement}, each pair contributes
total alignment at most~$1$, so the average
non-terminating alignment is at most~$1/2$. Therefore:
\[
(tm{-}1)\,\alpha(n) \le (m{-}1) +
\frac{(t{-}1)m}{2} = \frac{m(t{+}1)}{2} - 1.
\]
The ratio $(m(t{+}1)/2 - 1)/(tm - 1) \le (t{+}1)/(2t)$
holds for all $m \ge 1$ (the difference is
$(t{-}1)/(2t(tm{-}1)) \ge 0$).
\end{proof}

\begin{proposition}\label{prop:threshold}
In base~$10$, the rough part $t(n)$ is coprime to~$10$.
The coprime integers $t \ge 4$ with $\gcd(t, 10) = 1$
satisfy $t \ge 7$. For all such~$t$:
\[
\frac{t+1}{2t} \le \frac{8}{14} = \frac{4}{7}
< \frac{1}{\varphi}.
\]
\end{proposition}

\begin{proof}
The function $(t{+}1)/(2t) = 1/2 + 1/(2t)$ is
decreasing. Its maximum for $t \ge 7$ is at $t = 7$:
$8/14 = 4/7 \approx 0.5714$. Since
$1/\varphi \approx 0.6180$: $4/7 < 1/\varphi$.
\end{proof}

% ============================================================
\section{The Three-Tier Theorem}

\begin{theorem}[Three-tier classification]\label{thm:three-tier}
In base~$10$, the repetend alignment $\alpha(n) \ge 1/\varphi$
if and only if one of the following holds:
\begin{enumerate}
\item[\textup{(i)}] $n$ is $10$-smooth, or
\item[\textup{(ii)}] $n = 3m$ with $m$ a $10$-smooth integer and
      $m \ge 4$.
\end{enumerate}
\end{theorem}

\begin{proof}
\emph{Sufficiency.}
Case~(i): if $n$ is $10$-smooth, every fraction
$k/n$ terminates, so $\alpha(n) = 1 \ge 1/\varphi$.
Case~(ii): by~\cite[Theorem~2]{paper1},
$\alpha(3m) = (2m{-}1)/(3m{-}1) \ge 1/\varphi$ if and only if
$m \ge \varphi^{2} \approx 2.618$, hence $m \ge 3$. Since $m$ is
$10$-smooth and $m \ge 4$, the bound holds.

\medskip
\emph{Necessity.} Let $n$ have rough part $t$ and smooth
part $m$.

If $t \ge 7$: Lemma~\ref{lem:upper-bound} and
Proposition~\ref{prop:threshold} give
$\alpha(n) \le (t{+}1)/(2t) \le 4/7 < 1/\varphi$.

If $t = 3$ and $m \le 2$: direct computation gives
$\alpha(3) = 1/2$ and $\alpha(6) = 3/5$, both below
$1/\varphi$. (The value $m = 3$ is not $10$-smooth.)

If $t = 1$: $n$ is smooth, which is case~(i).

If $t = 3$ and $m \ge 4$ smooth: this is case~(ii).
\end{proof}

% ============================================================
\section{Remarks}

\subsection*{The gap}

The closest any Tier~3 integer comes to $1/\varphi$ is
$n = 6$ with $\alpha(6) = 3/5 = 0.600$. The gap is
$1/\varphi - 3/5 \approx 0.018$. For $t \ge 7$, the
nines complement bound gives $\alpha \le 4/7 \approx 0.571$,
leaving a gap of $1/\varphi - 4/7 \approx 0.047$.
The three-tier structure is robust.

\subsection*{The nines complement}

The proof rests on one observation: complementary fractions
$k/n$ and $(n{-}k)/n$ have digits summing to $b{-}1$ at
every position. For even base~$b$, this means they can
never simultaneously match a third fraction at the same
position. The average alignment of non-terminating
fractions with $1/n$ is therefore at most~$1/2$.

\subsection*{Why $t = 7$}

In base~$10$, the rough part $t$ must be coprime to~$10$.
The integers $4, 5, 6$ share a factor with~$10$, so
the smallest admissible $t \ge 4$ is $t = 7$. The bound
$(t{+}1)/(2t) = 4/7 < 1/\varphi$ closes the proof.

\subsection*{Other even bases}

The nines complement proof extends to even bases $b$
with $b \equiv 1 \pmod{3}$ and $3 \nmid b$: the
complement identity holds on the periodic denominator
(since $\gcd(t, b) = 1$), complementary digits are
distinct (since $b$ is even), and the smallest coprime
$t \ge 4$ satisfies $(t{+}1)/(2t) < 1/\varphi$.

\subsection*{Odd bases}

For odd bases, the three-tier classification requires
modification: the prime~$2$ is no longer $b$-smooth,
so $n = 2$ has $\alpha = 1$ but is not in Tier~1.
Additional families must be included. The complement
bound also requires care, since the middle digit
$(b{-}1)/2$ is an integer for odd~$b$. A counting
bound (bounding per-position matches by the bin size)
handles most cases, with finitely many $(b, t)$ pairs
verified directly.

% ============================================================
\begin{thebibliography}{9}

\bibitem{paper1}
A.~S.~Petty,
Why the golden ratio selects the prime three,
research note, 2020.

\bibitem{paper2}
A.~S.~Petty,
Digit-partitioning primes and the alignment formula,
research note, 2020.

\bibitem{paper3}
A.~S.~Petty,
The alignment limit for all primes,
research note, 2021.

\end{thebibliography}

\end{document}
